\section{Introduction}
\label{sec:intro}

Object detection, a foundational pillar of computer vision, aims to locate and identify objects within an image. Traditionally, object detection was operated within a closed-set paradigm~\cite{girshick2015fast, ren2015faster, redmon2016you, lin2017focal, tian2019fcos,zhou2019objects, zhu2020deformable, liu2022dabdetr, li2022dn, zhang2022dino, carion2020end}, wherein a predefined set of categories is known a prior, and the system is trained to recognize and detect objects from this set. Yet the ever-changing and unforeseeable nature of the real world demands a shift in object detection methodologies towards an open-set paradigm.

Open-set object detection represents a significant paradigm shift, transcending the limitations of closed-set detection by empowering models to identify objects beyond a predetermined set of categories. A prevalent approach is to use text prompts for open-vocabulary object detection~\cite{liu2023grounding, li2022grounded, ghiasi2022scaling, zhou2022detecting, gu2021open, minderer2022simple, kamath2021mdetr}. This approach typically involves distilling knowledge from language models like CLIP~\cite{clip} or BERT~\cite{devlin2018bert} to align textual descriptions with visual representations.                

\begin{figure}[t]
\centering
\includegraphics[width=0.95\linewidth]{images/intro.pdf}\vspace{-1mm}
\caption{Long-tailed curve of object frequency and the number of categories that can be detected. We suggest that the text prompt can cover the middle part of the long-tailed curve, while the visual prompt can cover the tail.}
\label{fig:intro1} 
\vspace{-2mm}
\end{figure}

While using text prompts has been predominantly favored in open-set detection for their capacity to abstractly describe objects, it still faces the following limitations. 
1) \textit{Long-tailed data shortage}. The training of text prompts necessitates modality alignment between visual representations, however, the scarcity of data for long-tailed objects may impair the learning efficiency. As depicted in ~\cref{fig:intro1}, the distribution of objects inherently follows a long-tail pattern, i.e., as the variety of detectable objects increases, the available data for these objects becomes increasingly scarce. This data scarcity may undermine the capacity of models to identify rare or novel objects. 2) \textit{Descriptive limitations}. Text prompts also fall short of accurately depicting objects that are hard to describe in language. For instance, as shown in~\cref{fig:intro1}, while a text prompt may effectively describe \texttt{ferris wheel}, it may struggle to accurately represent \texttt{the microorganisms in the microscope image} without biological knowledge.

 Conversely, visual prompts~\cite{xu2023multi, zou2023segment, li2023semantic, kirillov2023segment, li2023visual, jiang2023t} provide a more intuitive and direct method to represent objects by providing visual examples. For example, users can use points or boxes to mark the object for detection, even if they do not know what the object is. Additionally, visual prompts are not constrained by the need for cross-modal alignment, since they rely on visual similarity rather than linguistic correlation, enabling their application to novel objects that are not encountered during training.  

 Nonetheless, visual prompts also exhibit limitations, as they are less effective at capturing the general concept of objects compared to text prompts. For instance, the term \texttt{dog} as a text prompt broadly covers all dog varieties. In contrast, visual prompts, given the vast diversity in dog breeds, sizes, and colors, would necessitate a comprehensive image collection to visually convey the abstract notion of \texttt{dog}.

Recognizing the complementary strengths and weaknesses of both text and visual prompts, we introduce \ourmodel{}, a generic open-set object detection model that integrates both modalities. \ourmodel{} is built upon the DETR~\cite{carion2020end} architecture which is an end-to-end object detection model. It incorporates two parallel encoders to encode both text and visual prompts. For text prompts, we utilize the text encoder of CLIP~\cite{clip} to encode input text into text embedding. For visual prompts, we introduce a novel visual prompt encoder equipped with the deformable attention~\cite{zhu2020deformable} that can transform the input visual prompts (points or boxes) on a single image or across multiple images into visual prompt embeddings. To facilitate the collaboration of these two prompt modalities, we propose a contrastive learning~\cite{clip, hadsell2006dimensionality} module that can explicitly align text prompts and visual prompts. During the alignment, visual prompts can benefit from the generality and abstraction capabilities inherent in text prompts. Conversely, text prompts can enhance their descriptive capabilities by looking at various visual prompts. This iterative interaction allows both visual and text prompts to evolve continuously, thereby improving their ability for generic understanding within one model.

\ourmodel{} supports four unique workflows that can be applied to various scenarios: 1) \textit{interactive visual prompt workflow}, allowing users to specify the object to be detected by given visual example through box or point on the current image; 2) \textit{generic visual prompt workflow}, permitting users to define a specific object across multiple images through visual prompts, thereby creating a universal visual embedding applicable to other images; 3) \textit{text prompt workflow}, enabling users to employ descriptive text for open-vocabulary object detection; 4) \textit{mix prompt workflow}, which combines both text and visual prompts for joint inference. 


\ourmodel{} demonstrates strong object detection capabilities and achieves remarkable results on COCO~\cite{lin2014microsoft}, LVIS~\cite{gupta2019lvis}, ODinW~\cite{li2022elevater} and Roboflow100~\cite{ciaglia2022roboflow}, all under zero-shot setting. Through our analysis, we observe that text and visual prompts serve complementary roles, each excelling in scenarios where the other may not be as effective. Specifically, text prompts are particularly good at recognizing common objects, while visual prompts excel in rare objects or scenarios that may not be easily described through language. This complementary relationship enables the model to perform effectively across a wide range of scenarios. To summarize, our contributions are threefold:
\begin{itemize}
    \item We propose an open-set object detection model \ourmodel{} that unifies text and visual prompts within one framework, which demonstrates strong zero-shot capabilities across various scenarios.
    \item We propose a contrastive learning module to explicitly align text and visual prompts, which leads to mutual enhancement of these two modalities.
    \item Extensive experiments demonstrate the benefits of unifying text and visual prompts within one model. We also reveal that each type of prompt can cover different scenarios, which collectively show promise in advancing toward general object detection.
\end{itemize}